% Ad Familiares 15.9 --- headnote
% ad-familiares-15-09, summer 51 BC, in transit to Cilicia (companion to 15.7 and 15.8)

\textit{Cicero proconsul to M. Claudius Marcellus, consul of
51 BC, written from the same stage of the Cilician journey
that produced 15.7 and 15.8 --- the Perseus dateline reads
``at the same place and time as letters vii and viii.'' The
book-precision metadata in this corpus carries it generically
under 54 BC; the occasion belongs to summer or early autumn 51
BC, when news of the consular elections for 50 BC reached
Cicero on the road and he wrote in turn to M. Marcellus,
to his cousin C. Marcellus the consul-elect, and back to M.
Marcellus.}

\textit{Three sections, three notes. \textsection 1 is
congratulation: M. Marcellus has reaped the fruit of his own
distinguished consulship in seeing his cousin returned for
50 BC --- and Cicero, ``posted by you yourself to the ends
of the earth,'' raises him to heaven with the justest of
praises. The praise is balanced and intimate: Cicero loved
him in boyhood, and now hears from the best of men that the
two are alike --- ``vel me tui similem esse audio vel te mei,''
the chiasmic close of the section.}

\textit{\textsection 2 is the practical request that runs
through all the proconsular letters: do not let any time be
added to my term. The senate had limited Cicero's command to
one year; he is anxious already to be relieved. If Marcellus
will arrange a successor --- or simply prevent the fixed term
from being lengthened --- ``I shall reckon I have gained
everything through you.''}

\textit{\textsection 3 is the discreet aside, and the most
characteristic. News of the Parthians has reached him; he
will not yet write home officially; and for that reason ---
not even on the strength of friendship --- he will not even
write privately to a consul: ``lest, having written to a
consul, I should seem to have written officially.'' Public
and private overlap, in the consular's mailbag, in ways the
proconsul will not exploit.}
